
\subsection*{Contact feedback can arrive without being used effectively}

A further six-batch experiment separates access to consequences from the ability to exploit them in the same class of moving-node elastic host. Resource uptake at distal material emits a signal that travels at most one actual contact per update. In the opportunity run, the first uptake occurs at postclock 53 and the first proximal signal at 56, although the proximal material never takes up resource itself. Cutting only information transmission across a section in three actual parents preserves every recorded physical variable during 96 stopped-writing updates. The residual local signal is not erased. With writing enabled, the cut-minus-open uptake changes are tiny and of mixed sign; nearly full distal stores prevent this comparison from establishing a useful communication benefit (Supplementary Tables~\ref{tab:reach-means} and \ref{tab:reach-cut}).

Useful action consequences nevertheless exist. Three new bodies receive an actual 96-step resource-free prefix, then paid positive or negative proximal rest-length changes. Positive-minus-negative uptake is positive in all three bodies at resource position 5.4 (0.44420--0.45451) and negative at position 4.7 ($-0.42971$ to $-0.41614$). The received local signals retain the corresponding ordering. These are executed fixed-action comparisons, not counterfactual labels supplied to the learner, and they establish neither online credit assignment nor a unique causal role for one contact (Supplementary Table~\ref{tab:reach-actions}).

Providing reachable consequences does not make the tested ordinary rule reliably useful. Mean-age alignment lowers changed-environment uptake in all three histories compared with ordinary diffusion (means 0.75176 versus 1.23979). Holding actions for 96 instead of 24 updates improves that uptake in only one of three further histories; its higher terminal energy accompanies fewer paid parameter changes. Finally, three independent sparse-resource histories keep distal stores below capacity, yet diffusion yields less changed-environment uptake than no signal in all three, by 0.00542, 0.01881 and 0.00635 units. Instantaneous global feedback is not consistently better either. These unequal-information controls expose a failure of utilization under the tested rule, rather than a general absence of accessible consequences or a fair ranking of learning algorithms. Fixed material with a compensating readout exactly reproduces the first-batch physical states. Neither a new transmission-method advantage nor a broad impossibility of contact-mediated learning follows (Supplementary Tables~\ref{tab:reach-pairs}--\ref{tab:reach-jobs}).
